\section{蛋白质基础：序列、结构与功能}

\subsection{蛋白质是生物学的"执行器"}

蛋白质 (protein) 是一类极其多样的生物大分子，可以被视为\textbf{生物学的执行器} (biology's actuators)——它们驱动和调控细胞中几乎所有的功能。蛋白质的功能范围极广，包括但不限于：

\begin{itemize}
    \item \textbf{酶} (enzyme)：催化特定的化学反应，包括基因组水平的编辑（如 CRISPR 相关蛋白）
    \item \textbf{结构蛋白}：构成细胞骨架、细胞外基质等结构
    \item \textbf{信号蛋白}：传递细胞间和细胞内的信号（如钙离子结合蛋白）
    \item \textbf{治疗性蛋白}：如已获临床批准用于治疗白血病的蛋白质药物
    \item \textbf{转运蛋白}：在细胞内外运输分子
\end{itemize}

\begin{knowledgebox}{蛋白质作为治疗药物}
蛋白质不仅是天然的生物功能执行者，也可以被\textbf{工程化设计}为治疗性药物。例如，抗体药物（一种特殊的蛋白质）已被广泛用于癌症免疫治疗。Amini 在讲座中展示了一种已获临床批准、用于治疗特定类型白血病的蛋白质药物，说明蛋白质设计的临床转化价值已经得到验证。
\end{knowledgebox}

\subsection{蛋白质的三层表征层次}

蛋白质的生物学可以通过三个层次来理解，形成一个从线性编码到三维结构再到生物功能的层级关系：

\begin{table}[H]
\centering
\caption{蛋白质的三层表征层次}
\begin{tabular}{@{}lll@{}}
\toprule
\textbf{层次} & \textbf{描述} & \textbf{数据特征} \\
\midrule
序列 (Sequence) & 氨基酸的线性排列 & 离散符号序列，20 种氨基酸字母表 \\
结构 (Structure) & 三维空间构象 & 连续坐标，原子级几何信息 \\
功能 (Function) & 生物学活性与作用 & 多样化标注，难以量化 \\
\bottomrule
\end{tabular}
\end{table}

\textbf{序列层}：每个蛋白质由一条氨基酸 (amino acid) 序列定义。自然界存在 20 种标准氨基酸，每种可用单字母代码表示（如 A 代表丙氨酸 Alanine，G 代表甘氨酸 Glycine 等）。因此蛋白质序列可以被看作基于 20 个字符词汇表的一种"语言"，这为使用 NLP 技术处理蛋白质数据提供了天然的类比基础。

\textbf{结构层}：氨基酸序列通过复杂的物理化学过程折叠 (folding) 为特定的三维结构。结构决定了蛋白质如何与其他分子相互作用，是功能实现的物理基础。

\textbf{功能层}：蛋白质的生物学功能（如催化、结合、信号传导等）通常由其结构所决定，而结构又由序列编码。

\begin{importantbox}{中心法则的反向工程}
传统的蛋白质研究遵循 序列 $\rightarrow$ 结构 $\rightarrow$ 功能 的正向路径。然而，AI 驱动的蛋白质设计旨在\textbf{反转这一范式}：从我们期望的\textbf{功能}出发，逆向推导出实现该功能的\textbf{序列}。这正是生成式蛋白质设计的核心挑战：
$$\text{Function} \rightarrow \text{Sequence} \rightarrow \text{Structure (验证)}$$
\end{importantbox}

\subsection{现有的 AI 蛋白质工具}

在蛋白质 AI 领域，已有大量工作针对序列-结构-功能管线的各个环节：

\begin{itemize}
    \item \textbf{蛋白质语言模型} (Protein Language Models)：如 ESM 系列，对蛋白质序列进行表征学习，捕获进化和功能信息
    \item \textbf{结构生成模型}：如 RFdiffusion，能在结构空间中设计出自然界中不存在的新几何构型
    \item \textbf{序列到结构预测}：如 \textbf{AlphaFold}，给定氨基酸序列预测其三维结构——这是近年来计算生物学领域最重大的突破之一
\end{itemize}

\begin{warningbox}{结构数据与序列数据的规模鸿沟}
蛋白质领域存在严重的\textbf{数据不对称}：
\begin{itemize}
    \item \textbf{序列数据}：数亿乃至数十亿条独特蛋白质序列（来自基因组测序）
    \item \textbf{结构数据}：仅约 30 万个实验解析的蛋白质结构（来自 PDB 数据库）
\end{itemize}
虽然 AlphaFold 可以预测结构，但那是\textbf{计算预测}而非实验测量。结构数据库中的蛋白质倾向于球状 (globular)、可溶性蛋白，富含 $\alpha$-helix 结构元素，导致基于结构数据训练的模型存在系统性偏差。
\end{warningbox}

\subsection{本章小结}

蛋白质是连接序列编码与生物功能的核心生物大分子。理解蛋白质的序列-结构-功能三层层级关系，是构建 AI 蛋白质设计工具的基础。现有工具已能分别处理管线的各个环节，但实现从功能到序列的反向设计仍是开放挑战。数据层面，序列数据的丰富性远超结构数据，这一不对称性深刻影响着模型设计策略的选择。


